\documentclass[10pt]{article}
\usepackage[a4paper,margin=23mm]{geometry}
\usepackage{amsmath,amssymb}
\usepackage[hidelinks]{hyperref}
\setlength{\emergencystretch}{3em}
\title{Conjecture 00000001546: a polynomial growth obstruction}
\author{lizaixi01, with OpenAI Codex assistance}
\date{23 October 2026}
\begin{document}
\maketitle
\section{Two incompatible necessary clauses}
Both source languages assert that the maximum bisector-line count
has upper growth $n e^{O(\sqrt{\log n})}$ and that no polynomial
upper bound exists. These clauses are incompatible for \emph{any}
candidate counting function. We use the ordinary meaning of a
polynomial upper bound: $O(n^d)$ for some fixed nonnegative integer
$d$. In fact, the asserted upper growth implies $O(n^2)$.

This does not assign a value to a geometric counting function or
replace bisector lines by an artificial model. It refutes the
necessary numerical conjunction regardless of what the actual
count is. We do not determine the correct maximal halving-line
count, resolve variants of the source's geometric definition, or
deny the compatible lower growth $n e^{\Omega(\sqrt{\log n})}$.

\section{An explicit analytic bound}
Let $x\ge1$ and $C\in\mathbb R$. Set $u=\sqrt{\log x}$.
Since $(u-C/2)^2\ge0$ and $u^2=\log x$,
\[
 C\sqrt{\log x}\le\log x+\frac{C^2}{4}.
\]
Exponentiating and multiplying by $x\ge0$ gives the exact bound
\[
 x e^{C\sqrt{\log x}}\le e^{C^2/4}x^2. \tag{1}
\]
This holds for every real $C$, with an explicit constant and the
fixed cutoff $x\ge1$. Hence every function that is
$O(n e^{C\sqrt{\log n}})$ for some $C$ is $O(n^2)$.

The proof also treats the $O$ \emph{inside} the exponential
directly. If $h(n)=O(\sqrt{\log n})$, its genuine eventual norm
bound gives $h(n)\le C\sqrt{\log n}$ beyond a cutoff. Apply (1)
there to obtain
\[
 n e^{h(n)}=O(n^2).
\]
Any eventually equal function has the same bound. For a one-sided
counting bound $f(n)\le n e^{C\sqrt{\log n}}$, Lean explicitly
requires $f(n)\ge0$; the natural-valued specialization supplies
this automatically. No negative-function loophole is used.

\section{Formal conclusion and verification}
Mathlib's actual \texttt{Asymptotics.IsBigO} at \texttt{atTop}
encodes all bound constants and cutoffs. The final statements
quantify over every candidate real-valued function, and separately
every natural-valued count, and deny the conjunction of the printed
upper-growth condition with absence of every $O(n^d)$. The witness
polynomial exponent is the actual integer $d=2$. Real logarithm,
square root, exponential, eventual equality and Big-$O$
transitivity are all proved with their actual analytic meanings.

Lean 4.33.0 and Mathlib commit
\begin{center}\small\texttt{db584cd6d46c92f209a44c0f1c829460d327499d}\end{center}
are pinned. A new root build, all eleven theorem axiom audits and
full official \texttt{leanchecker --fresh Main} replay are included.
No \texttt{sorry}, \texttt{native\_decide} or additional axiom is
used. OpenAI Codex assisted, with a separate supervising Codex
agent reviewing source, proof, report and rendered PDF. This is
AI review, not human certification. Organizer acceptance is pending.
\end{document}
